\documentclass[11pt]{article}
% ===========================================================================
%  Shared preamble for the AP Calculus AB exam-preparation worksheets.
%  Input right after \documentclass. Files beginning with "_" are not compiled
%  on their own. Navy and gold, to match the course's lesson pages.
% ===========================================================================
\usepackage[letterpaper,top=2.3cm,bottom=2.2cm,left=1.7cm,right=1.7cm,headheight=28pt,footskip=30pt]{geometry}
\usepackage{amsmath,amssymb}
\usepackage[T1]{fontenc}
\usepackage[utf8]{inputenc}
\usepackage{lmodern}
\usepackage{xcolor}
\usepackage[most]{tcolorbox}
\usepackage{enumitem}
\usepackage{fancyhdr}
\usepackage{array}
\usepackage{tikz}
\usepackage{pgfplots}
\pgfplotsset{compat=1.18}

\definecolor{navy}{HTML}{161E54}
\definecolor{ink}{HTML}{27357E}
\definecolor{gold}{HTML}{B98A00}
\definecolor{goldbg}{HTML}{FFFDF3}
\definecolor{inkbg}{HTML}{F6F7FD}
\definecolor{rule}{HTML}{C9CFF0}
\definecolor{okgreen}{HTML}{1B7A44}
\definecolor{okbg}{HTML}{F1F8F3}

\pagestyle{fancy}
\fancyhf{}
\fancyhead[L]{\footnotesize Name: \underline{\hspace{4cm}}}
\fancyhead[C]{\textbf{\color{navy}Integration Academy}\\[-2pt]{\scriptsize Calculus AB \textperiodcentered{} AP\textsuperscript{\textregistered} Exam Preparation}}
\fancyfoot[L]{\tiny\color{black!55}\parbox[t]{0.78\textwidth}{AP\textsuperscript{\textregistered} is a registered trademark of the College Board, which was not involved in the production of, and does not endorse, this material. \textcopyright{} Integration Academy.}}
\fancyfoot[R]{\scriptsize\thepage}
\renewcommand{\headrulewidth}{1.2pt}
\renewcommand{\headrule}{\hbox to\headwidth{\color{gold}\leaders\hrule height \headrulewidth\hfill}}
\renewcommand{\footrulewidth}{0pt}
\setlength{\parindent}{0pt}
\setlength{\parskip}{4pt}

% Title block: \wstitle{code}{title}{unit line}
\newcommand{\wstitle}[3]{\begin{center}{\small\bfseries\color{gold}WORKSHEET #1}\\[3pt]{\LARGE\bfseries\color{navy}#2}\\[3pt]{\small\color{black!60}#3}\end{center}\vspace{2pt}}
% Section heading
\newcommand{\sech}[1]{\par\vspace{10pt}{\large\bfseries\color{navy}#1}\par\nobreak\vspace{2pt}{\color{rule}\hrule height0.8pt}\vspace{7pt}}
% A level tag printed beside a group of questions
\newcommand{\level}[1]{\par\vspace{6pt}{\footnotesize\bfseries\color{gold}\MakeUppercase{#1}}\par\nobreak\vspace{2pt}}

\newtcolorbox{ideabox}{colframe=ink,colback=inkbg,coltitle=white,colbacktitle=ink,fonttitle=\bfseries\small,title={Key ideas},arc=2pt,boxrule=0.8pt,left=8pt,right=8pt}
\newtcolorbox{example}[1]{colframe=rule,colback=white,coltitle=navy,colbacktitle=inkbg,fonttitle=\bfseries\small,title={#1},arc=2pt,boxrule=0.8pt,borderline west={3pt}{0pt}{ink},left=10pt,right=8pt}
\newtcolorbox{frbox}[1]{colframe=gold,colback=goldbg,coltitle=white,colbacktitle=gold,fonttitle=\bfseries\small,title={#1},arc=2pt,boxrule=0.8pt,breakable,left=8pt,right=8pt}
\newtcolorbox{solbox}{colframe=okgreen,colback=okbg,coltitle=white,colbacktitle=okgreen,fonttitle=\bfseries\small,title={Worked solutions},arc=2pt,boxrule=0.8pt,breakable,left=8pt,right=8pt}
\newtcolorbox{scorebox}[1]{colframe=navy,colback=white,coltitle=white,colbacktitle=navy,fonttitle=\bfseries\small,title={#1},arc=2pt,boxrule=0.8pt,breakable,left=8pt,right=8pt}

\newcommand{\soln}{\par\smallskip\textbf{\color{ink}Solution.}\ }
\newcommand{\result}[1]{\par\smallskip\textbf{#1}}
% blank working space of a given height
\newcommand{\work}[1]{\par\vspace{#1}}
% a point value in the margin of a scoring guide
\newcommand{\pts}[1]{\hfill{\footnotesize\bfseries\color{gold}[#1]}}

% Coordinate grid for graph questions: \begin{apgraph}{xmin}{xmax}{ymin}{ymax} ... \end{apgraph}
\newenvironment{apgraph}[4]{\begin{center}\begin{tikzpicture}\begin{axis}[
  axis lines=middle, xmin=#1, xmax=#2, ymin=#3, ymax=#4,
  grid=both, grid style={black!12}, xtick distance=1, ytick distance=1,
  ticklabel style={font=\scriptsize}, xlabel={$x$}, ylabel={$y$},
  width=10cm, height=7.2cm, samples=80, clip=false,
  every axis plot/.append style={line width=1.1pt, color=ink}]}{\end{axis}\end{tikzpicture}\end{center}}
\newcommand{\closeddot}[2]{\addplot[only marks, mark=*, mark size=2.2pt, color=ink] coordinates {(#1,#2)};}
\newcommand{\opendot}[2]{\addplot[only marks, mark=*, mark size=2.2pt, color=ink, mark options={fill=white}] coordinates {(#1,#2)};}

% Multiple-choice options: four across, or stacked when the options are long.
\newcommand{\opts}[4]{\par\vspace{3pt}\begin{tabular}{@{}p{0.235\linewidth}p{0.235\linewidth}p{0.235\linewidth}p{0.235\linewidth}@{}}(A)\ #1 & (B)\ #2 & (C)\ #3 & (D)\ #4\end{tabular}\par\vspace{9pt}}
\newcommand{\optsv}[4]{\par\vspace{3pt}\begin{tabular}{@{}l@{\ \ }p{0.9\linewidth}@{}}(A) & #1\\(B) & #2\\(C) & #3\\(D) & #4\end{tabular}\par\vspace{9pt}}

\fancyhead[R]{\footnotesize\color{black!60}Worksheet 8.1}
\begin{document}

\wstitle{8.1}{Average Value of a Function}{Unit 8: Applications of Integration}

\begin{ideabox}
\begin{itemize}[leftmargin=*,itemsep=2pt,topsep=2pt]
  \item The \textbf{average value} of $f$ on $[a,b]$ is $\dfrac{1}{b-a}\displaystyle\int_{a}^{b}f(x)\,dx$.
  \item It is the height of the rectangle on $[a,b]$ with the same area as the region under the graph.
  \item If $f$ is continuous on $[a,b]$, then $f(c)$ equals the average value for at least one $c$ in the interval.
  \item Do not confuse it with the \textbf{average rate of change}, $\dfrac{f(b)-f(a)}{b-a}$, which is the average value of $f'$.
\end{itemize}
\end{ideabox}

\sech{Worked examples}
\begin{example}{Example 1 \textperiodcentered{} A polynomial}
Find the average value of $f(x)=3x^2-2x$ on $[1,4]$.\soln $\displaystyle\int_{1}^{4}(3x^2-2x)\,dx=\Big[x^3-x^2\Big]_1^4=48$.\result{$\dfrac{48}{3}=16$.}
\end{example}

\begin{example}{Example 2 \textperiodcentered{} Where the average is reached}
Find the value of $c$ in $[0,9]$ at which $f(x)=\sqrt{x}$ equals its average value on $[0,9]$.\soln $\displaystyle\int_{0}^{9}\sqrt x\,dx=18$, so the average is 2. Solve $\sqrt c=2$.\result{$c=4$.}
\end{example}

\begin{example}{Example 3 \textperiodcentered{} An exponential}
Find the average value of $f(x)=e^{2x}$ on $[0,1]$.\soln $\dfrac{1}{1}\Big[\tfrac12e^{2x}\Big]_0^1$.\result{$\dfrac{e^2-1}{2}$.}
\end{example}

\newpage
\sech{Practice}

\level{Core}
\begin{enumerate}[leftmargin=*,itemsep=2pt]
  \item Find the average value of $f(x)=6x$ on $[0,4]$.\work{2.3cm}
  \item Find the average value of $f(x)=x^3$ on $[0,2]$.\work{2.3cm}
  \item Find the average value of $f(x)=\dfrac{1}{x^2}$ on $[1,4]$.\work{2.3cm}
  \item Given $\displaystyle\int_{2}^{7}f(x)\,dx=30$, find the average value of $f$ on $[2,7]$.\work{2.3cm}
\end{enumerate}

\level{Exam level}
\begin{enumerate}[leftmargin=*,itemsep=2pt,start=5]
  \item Find the average value of $f(x)=x\sqrt{x^2+9}$ on $[0,4]$.\work{2.8cm}
  \item The average value of $f$ on $[1,5]$ is 7. Find $\displaystyle\int_{1}^{5}\big(f(x)+2\big)\,dx$.\work{2.8cm}
  \item Find the value of $c$ in $[0,6]$ at which $f(x)=x^2+2$ equals its average value on $[0,6]$.\work{2.8cm}
  \item The table gives the temperature $H(t)$, in $^\circ$C, at selected times $t$ hours. Use a trapezoidal sum with three subintervals to estimate the average temperature for $0\le t\le10$.\begin{center}\renewcommand{\arraystretch}{1.25}\begin{tabular}{|c|c|c|c|c|}\hline $t$ & $0$ & $2$ & $6$ & $10$ \\ \hline $H(t)$ & $12$ & $18$ & $20$ & $14$ \\ \hline\end{tabular}\end{center}\work{2.8cm}
\end{enumerate}

\begin{frbox}{Question 9 \textperiodcentered{} Free-response practice}
\emph{Calculator permitted.} The temperature in a greenhouse is $W(t)=20+15\sin\!\left(\dfrac{\pi t}{12}\right)$ degrees Celsius, for $0\le t\le12$, where $t$ is in hours.
\begin{enumerate}[label=(\alph*),leftmargin=*,itemsep=2pt]
  \item Find the average temperature for $0\le t\le12$.\work{2.6cm}
  \item Find the average rate of change of the temperature for $0\le t\le6$. Include units.\work{2.6cm}
  \item Find the times at which the temperature equals its average value.\work{2.6cm}
\end{enumerate}
\end{frbox}

\newpage
\begin{solbox}
\begin{enumerate}[leftmargin=*,itemsep=4pt]
  \item $\dfrac14(48)=12$.
  \item $\dfrac12(4)=2$.
  \item $\dfrac13\left(1-\dfrac14\right)=\dfrac14$.
  \item $\dfrac{30}{5}=6$.
  \item $\dfrac14\Big[\tfrac13(x^2+9)^{3/2}\Big]_0^4=\dfrac14\cdot\dfrac{125-27}{3}=\dfrac{49}{6}$.
  \item $\displaystyle\int_{1}^{5}f(x)\,dx=28$, and $\displaystyle\int_{1}^{5}2\,dx=8$: total 36.
  \item Average $=\dfrac16(72+12)=14$; $c^2+2=14$ gives $c=2\sqrt3$.
  \item $\dfrac{2(12+18)}{2}+\dfrac{4(18+20)}{2}+\dfrac{4(20+14)}{2}=174$; average $\approx17.4$ $^\circ$C.
  \item \begin{enumerate}[label=(\alph*),leftmargin=*,itemsep=3pt]
    \item $\dfrac{1}{12}\displaystyle\int_{0}^{12}W(t)\,dt=20+\dfrac{30}{\pi}\approx29.549$ $^\circ$C.
    \item $\dfrac{W(6)-W(0)}{6}=\dfrac{35-20}{6}=2.5$ $^\circ$C per hour.
    \item $\sin\!\left(\dfrac{\pi t}{12}\right)=\dfrac2\pi$: $t\approx2.636$ and $t\approx9.364$ hours.
  \end{enumerate}
\end{enumerate}
\end{solbox}

\end{document}
